% =============================================================================
%  Section: Event and Track Selection
%  H1 Charged-Hadron SIDIS High-Q^2 Analysis Note
% =============================================================================
\section{Event and Track Selection}
\label{sec:selection}

The analysis uses the standard H1 high-$\Qsq$ neutral-current DIS event
selection followed by a reconstructed-particle selection for the
charged-hadron measurement.  The event preselection defines the accepted event
sample that is passed to the skimmed analysis TTree and, for simulated samples, keeps additional reconstructed and generator-level
information needed to model detector migrations.  The final analysis selection
is applied identically to the measured reconstructed-particle sample and to the
reconstructed side of the response-matrix construction.

\subsection{DIS Event Preselection}
\label{sec:selection:preselection}

The DIS event preselection requires a well reconstructed scattered-electron
candidate, a good primary vertex, and rejection of non-$ep$ and QED-Compton-like
backgrounds.  The requirements are:
\begin{itemize}
  \item The good-run list and detector-status requirements are applied;
  \item The S67 trigger (electron in the LAr calorimeter) is required;
  \item The scattered-electron energy satisfies $E_e \ge 11.0~\GeV$, and the
        scattered-electron polar angle satisfies $\theta_e < 2.7$ rad;
  \item The event vertex is available from the standard neutral-current vertex
        reconstruction: for forward electrons with $\theta_e<30^\circ$, the CJC
        fallback vertex is also accepted; in either case
        $|z_{\mathrm{vtx}}| \le 35.0~\mathrm{cm}$ is required;
  \item The loose cut envelope used in generation of the skimmed analysis TTree is
        $100 < Q^2_{e\Sigma} \le 20000~\GeVsq$,
        $0.05 < y_{e\Sigma} < 0.9$, $30 < E-p_z < 75~\GeV$, and
        $y_e < 0.94$;
  \item When bunch-type information is available, the colliding-bunch (i.e. not the pilot bunches) category 
        is required.
\end{itemize}

The scattered electron is required to fall in well instrumented regions of the
liquid-argon calorimeter.  The longitudinal impact coordinate must satisfy
\begin{equation}
  -190.0 < z_{\mathrm{impact}} < 15.0
  \quad\mathrm{or}\quad
  z_{\mathrm{impact}} > 25.0,
\end{equation}
and the azimuthal position within a calorimeter octant must avoid the octant
cracks according to
\begin{equation}
  2^\circ < \phi_{\mathrm{octant}} < 43^\circ
  \quad\mathrm{or}\quad
  z_{\mathrm{impact}} > 300.0.
\end{equation}
In addition, a run-dependent electromagnetic-calorimeter fiducial veto is
applied.  The veto is defined by the run number, the electron azimuthal position
at the liquid-argon calorimeter edge, the electron longitudinal impact
coordinate, and the electron energy.  It removes regions where the electron
trigger efficiency is below 100\%.

% Provenance for the non-ep veto fractions quoted here:
% /data/dust/group/h1/klesthen/BaryonAsym/H1ChargedSidis_HighQ2/output_v6/studies/bgfinder_ptratio_cutflow_2026-10-06/merged/summary_bgfinder_cutflow.json
A small fraction of the recorded events do not originate from $ep$ collisions,
but from interactions of the proton beam with residual gas in the beam pipe,
from beam-halo muons, or from cosmic rays traversing the detector.  Such
non-$ep$ backgrounds can mimic DIS events, for example when an unrelated
energy deposit is reconstructed as the scattered electron, but they are
characterised by event times inconsistent with the nominal beam crossing
and by a hadronic system that does not balance the transverse momentum of
the apparent electron candidate.  These events are rejected with the standard H1
timing vetoes: events flagged by the background finders
(flags 0, 1, 5, 6, 7), which combine timing information from the central tracker and the LAr calorimeter with topological information, are removed in
combination with the hadronic-to-electromagnetic transverse-momentum balance
$R=p_T^{\mathrm{had}}/p_T^{e}$.  

A single finder flag is treated as weak
evidence, since the finders occasionally also flag genuine $ep$ events, and
leads to rejection only when corroborated by a significant momentum
imbalance, $R<0.1$ for finder 7 and $R<0.5$ for finders 5 or 6, whereas two
simultaneous finder flags are treated as sufficient evidence and lead to
rejection for events with $R>0.1$.  The veto rejects 0.29\% of the events
passing all preceding selection requirements, while 1.61\% of these events
have $R<0.5$ (0.06\% have $R<0.1$) and are retained in the absence of a
corroborating finder flag, so the timing and kinematic criteria act jointly
rather than as independent requirements.  The anti-QED-
Compton veto rejects events with two electromagnetic particles when the two
leading electron candidates each have energy above $4~\GeV$, their summed energy
is above $18~\GeV$, and the two candidates are nearly back-to-back in azimuth,
$-\cos(|\phi_1-\phi_2|)>0.95$.

\subsection{Final Reconstructed-Event Selection}
\label{sec:selection:final-events}

The final reconstructed-event selection tightens the DIS phase space used for
the measurement:
\begin{itemize}
  \item $150 < Q^2 < 20000~\GeVsq$;
  \item $0.1 < y < 0.7$;
  \item $E-p_z > 45~\GeV$;
  \item $|z_{\mathrm{vtx}}| < 35~\mathrm{cm}$;
\end{itemize}
The same reconstructed-event requirements are used for the response-matrix
construction and for the measured event sample.

\subsection{Final Reconstructed-Track Selection}
\label{sec:selection:final-tracks}

The measured charged-hadron sample consists of reconstructed charged tracks
passing the track-quality, topology, and analysis-bin requirements.  Accepted tracks are counted only as inclusive $\hplus$ or $\hminus$ candidates according to the sign of their electric charge.
The truth-level charged-hadron convention used for acceptance and migration
modeling is documented in Sec.~\ref{sec:species_definition}.

The final reconstructed-track cuts are:
\begin{itemize}
  \item Lab-frame transverse momentum $p_T>0.15~\GeV$ and pseudorapidity $|\eta|<1.6$;
  \item At least 5 hits in the CJC
  \item Radial track length at least 5 cm;
  \item $|\mathrm{DCA}'|\sin\theta < 2$ cm,
  \item The track is labeled as originating from the primary vertex in H1OO;
  \item The reconstructed track is in the current hemisphere, i.e. $x_F<0$ in the H1 convention;
  \item The reconstructed $z$ and $\gamma^*p$-frame $\pTh$ are finite and inside
        the analysis binning.
   \item The opening angle between the track and the scattered electron is more than 0.2 radians.
\end{itemize}
The analysis binning uses
$z=[0.0,0.05,0.1,0.2,0.3,0.5,0.7,1.0]$ and reconstructed
$\pTh=[0.0,0.25,0.5,0.75,1.0,1.5,2.0,3.0,4.0,6.0,8.0]~\GeV$ edges.  The first
physical $z$ bin is $0.0 \le z < 0.05$.

A reconstructed-side opening-angle requirement suppresses tracks close to the
scattered-electron direction, which can come from final-state radiated photon conversions or mis-identification of the scattered electron track.  For each reconstructed track, the opening angle
to the scattered-electron three-vector, $\alpha$, is calculated in the lab frame. Accepted tracks must satisfy
$\alpha(\mathrm{track},e') \ge 0.2~\mathrm{rad}$.  This requirement is applied
only to reconstructed tracks; the corresponding generator-level loss of true charged hadrons is 0.106\% for RAPGAP and 0.114\% for DJANGOH.
